% =========================================================
% One-Sided Limits --- Notes
% Chapter: functions
% Section: one-sided-limits
% Volume III $\cdot$ Analysis
% =========================================================

\subsection*{One-Sided Limits}
\label{sec:one-sided-limits}
% ---------------------------------------------------------
% Definition --- One-Sided Limits
% ---------------------------------------------------------

\begin{definitionbox}{Definition (Right-Hand Limit)}
\begin{definition}[Right-Hand Limit]
\label{def:right-hand-limit}
Let $A\subseteq\mathbb{R}$, let $f:A\to\mathbb{R}$, let $c\in\mathbb{R}$,
and suppose $c$ is a cluster point of $A\cap(c,\infty)$. A real number
$L$ is the \textbf{right-hand limit} of $f$ at $c$ if
\[
  \forall\varepsilon>0\;\exists\delta>0\;\forall x\in A\;
  0<x-c<\delta\Rightarrow |f(x)-L|<\varepsilon.
\]
\end{definition}
\end{definitionbox}

\begin{remark*}[Standard quantified statement]
\[
  \forall\varepsilon>0\;\exists\delta>0\;\forall x\in A\;
  c<x<c+\delta\Rightarrow |f(x)-L|<\varepsilon.
\]
\end{remark*}

\begin{remark*}[Predicate reading]
\[
  \operatorname{RightLimit}(f,c,L)
  \;\coloneqq\;
  \forall\varepsilon>0\;\exists\delta>0\;\forall x\in A\;
  0<x-c<\delta\Rightarrow |f(x)-L|<\varepsilon.
\]
\end{remark*}

\begin{remark*}[Negated quantified statement]
\[
  \exists\varepsilon_0>0\;\forall\delta>0\;\exists x\in A\;
  0<x-c<\delta\;\wedge\; |f(x)-L|\geq\varepsilon_0.
\]
\end{remark*}

\begin{remark*}[Negation predicate reading]
\[
  \neg\,\operatorname{RightLimit}(f,c,L).
\]
\end{remark*}

\begin{remark*}[Failure modes]
A proposed right-hand limit fails when some positive output tolerance is
violated arbitrarily close to $c$ from the right.
\end{remark*}

\begin{remark*}[Failure modes]
\[
  \exists\varepsilon_0>0\;\forall\delta>0\;\exists x\in A\;
  0<x-c<\delta\;\wedge\; |f(x)-L|\geq\varepsilon_0.
\]
\end{remark*}

\begin{remark*}[Interpretation]
Only inputs to the right of $c$ are tested. The value of $f$ at $c$, if
it is defined, is irrelevant to the right-hand limit.
\end{remark*}

\begin{dependencies}
\begin{itemize}
  \item \hyperref[def:limit-function]{Limit of a Function}
  \item \hyperref[def:cluster-point-r]{Cluster Point}
\end{itemize}
\end{dependencies}

\begin{definitionbox}{Definition (Left-Hand Limit)}
\begin{definition}[Left-Hand Limit]
\label{def:left-hand-limit}
Let $A\subseteq\mathbb{R}$, let $f:A\to\mathbb{R}$, let $c\in\mathbb{R}$,
and suppose $c$ is a cluster point of $A\cap(-\infty,c)$. A real number
$L$ is the \textbf{left-hand limit} of $f$ at $c$ if
\[
  \forall\varepsilon>0\;\exists\delta>0\;\forall x\in A\;
  0<c-x<\delta\Rightarrow |f(x)-L|<\varepsilon.
\]
\end{definition}
\end{definitionbox}

\begin{remark*}[Standard quantified statement]
\[
  \forall\varepsilon>0\;\exists\delta>0\;\forall x\in A\;
  c-\delta<x<c\Rightarrow |f(x)-L|<\varepsilon.
\]
\end{remark*}

\begin{remark*}[Predicate reading]
\[
  \operatorname{LeftLimit}(f,c,L)
  \;\coloneqq\;
  \forall\varepsilon>0\;\exists\delta>0\;\forall x\in A\;
  0<c-x<\delta\Rightarrow |f(x)-L|<\varepsilon.
\]
\end{remark*}

\begin{remark*}[Negated quantified statement]
\[
  \exists\varepsilon_0>0\;\forall\delta>0\;\exists x\in A\;
  0<c-x<\delta\;\wedge\; |f(x)-L|\geq\varepsilon_0.
\]
\end{remark*}

\begin{remark*}[Negation predicate reading]
\[
  \neg\,\operatorname{LeftLimit}(f,c,L).
\]
\end{remark*}

\begin{remark*}[Failure modes]
A proposed left-hand limit fails when some positive output tolerance is
violated arbitrarily close to $c$ from the left.
\end{remark*}

\begin{remark*}[Failure modes]
\[
  \exists\varepsilon_0>0\;\forall\delta>0\;\exists x\in A\;
  0<c-x<\delta\;\wedge\; |f(x)-L|\geq\varepsilon_0.
\]
\end{remark*}

\begin{remark*}[Interpretation]
Only inputs to the left of $c$ are tested. This is the mirror image of
the right-hand limit definition.
\end{remark*}

\begin{dependencies}
\begin{itemize}
  \item \hyperref[def:limit-function]{Limit of a Function}
  \item \hyperref[def:cluster-point-r]{Cluster Point}
\end{itemize}
\end{dependencies}

% ---------------------------------------------------------
% Theorem --- Sequential Criterion for Right-Hand Limits
% ---------------------------------------------------------

\begin{theorembox}{Theorem}
\begin{theorem}[Sequential Criterion for Right-Hand Limits]
\label{thm:sequential-criterion-right-hand-limit}
Let $A\subseteq\mathbb{R}$, let $f:A\to\mathbb{R}$, and suppose $c$ is a
cluster point of $A\cap(c,\infty)$. Then $\lim_{x\to c^+}f(x)=L$ if and
only if for every sequence $(x_n)\subseteq A\cap(c,\infty)$ with
$x_n\to c$, one has $f(x_n)\to L$.

\smallskip\noindent
\hyperref[prf:sequential-criterion-right-hand-limit]{\textit{Go to proof.}}
\end{theorem}
\end{theorembox}

\begin{remark*}[Standard quantified statement]
\[
\begin{aligned}
&\Bigl[\forall\varepsilon>0\;\exists\delta>0\;\forall x\in A\;
  0<x-c<\delta\Rightarrow |f(x)-L|<\varepsilon\Bigr] \\
&\quad\Longleftrightarrow
\Bigl[\forall (x_n)\subseteq A\cap(c,\infty)\;
  x_n\to c\Rightarrow f(x_n)\to L\Bigr].
\end{aligned}
\]
\end{remark*}

\begin{remark*}[Predicate reading]
\[
  \operatorname{RightLimit}(f,c,L)
  \;\Longleftrightarrow\;
  \forall (x_n)\subseteq A\cap(c,\infty)\;
  x_n\to c\Rightarrow f(x_n)\to L.
\]
\end{remark*}

\begin{remark*}[Interpretation]
Right-hand limits can be tested by sequences approaching $c$ from above.
Every such sequence must force the corresponding output sequence toward
$L$.
\end{remark*}

\begin{dependencies}
\begin{itemize}
  \item \hyperref[def:right-hand-limit]{Right-Hand Limit}
  \item \hyperref[def:limit-function-sequential]{Sequential Definition of Limit}
\end{itemize}
\end{dependencies}

% ---------------------------------------------------------
% Corollary --- Sequential Criterion for Left-Hand Limits
% ---------------------------------------------------------

\begin{corollarybox}{Corollary}
\begin{corollary}[Sequential Criterion for Left-Hand Limits]
\label{cor:sequential-criterion-left-hand-limit}
Let $A\subseteq\mathbb{R}$, let $f:A\to\mathbb{R}$, and suppose $c$ is a
cluster point of $A\cap(-\infty,c)$. Then $\lim_{x\to c^-}f(x)=L$ if and
only if for every sequence $(x_n)\subseteq A\cap(-\infty,c)$ with
$x_n\to c$, one has $f(x_n)\to L$.

\smallskip\noindent
\hyperref[prf:sequential-criterion-left-hand-limit]{\textit{Go to proof.}}
\end{corollary}
\end{corollarybox}

\begin{remark*}[Standard quantified statement]
\[
\begin{aligned}
&\Bigl[\forall\varepsilon>0\;\exists\delta>0\;\forall x\in A\;
  0<c-x<\delta\Rightarrow |f(x)-L|<\varepsilon\Bigr] \\
&\quad\Longleftrightarrow
\Bigl[\forall (x_n)\subseteq A\cap(-\infty,c)\;
  x_n\to c\Rightarrow f(x_n)\to L\Bigr].
\end{aligned}
\]
\end{remark*}

\begin{remark*}[Predicate reading]
\[
  \operatorname{LeftLimit}(f,c,L)
  \;\Longleftrightarrow\;
  \forall (x_n)\subseteq A\cap(-\infty,c)\;
  x_n\to c\Rightarrow f(x_n)\to L.
\]
\end{remark*}

\begin{remark*}[Interpretation]
Left-hand limits can be tested by sequences approaching $c$ from below.
The criterion is the left-sided analogue of the right-hand sequential
criterion.
\end{remark*}

\begin{dependencies}
\begin{itemize}
  \item \hyperref[def:left-hand-limit]{Left-Hand Limit}
  \item \hyperref[def:limit-function-sequential]{Sequential Definition of Limit}
\end{itemize}
\end{dependencies}

% ---------------------------------------------------------
% Theorem --- Two-Sided Limit via One-Sided Limits
% ---------------------------------------------------------

\begin{theorembox}{Theorem}
\begin{theorem}[Two-Sided Limit via One-Sided Limits]
\label{thm:two-sided-limit-iff-matching-one-sided-limits}
Let $A\subseteq\mathbb{R}$ and let $f:A\to\mathbb{R}$. If $c$ is a
cluster point of both $A\cap(c,\infty)$ and $A\cap(-\infty,c)$, then
$\lim_{x\to c}f(x)=L$ if and only if both $\lim_{x\to c^+}f(x)=L$ and
$\lim_{x\to c^-}f(x)=L$.

\smallskip\noindent
\hyperref[prf:two-sided-limit-iff-matching-one-sided-limits]{\textit{Go to proof.}}
\end{theorem}
\end{theorembox}

\begin{remark*}[Standard quantified statement]
\[
\begin{aligned}
&\Bigl[\forall\varepsilon>0\;\exists\delta>0\;\forall x\in A\;
  0<|x-c|<\delta\Rightarrow |f(x)-L|<\varepsilon\Bigr] \\
&\quad\Longleftrightarrow
\Bigl[\forall\varepsilon>0\;\exists\delta_+>0\;\forall x\in A\;
  0<x-c<\delta_+\Rightarrow |f(x)-L|<\varepsilon\Bigr] \\
&\qquad\wedge
\Bigl[\forall\varepsilon>0\;\exists\delta_->0\;\forall x\in A\;
  0<c-x<\delta_-\Rightarrow |f(x)-L|<\varepsilon\Bigr].
\end{aligned}
\]
\end{remark*}

\begin{remark*}[Predicate reading]
\[
  \operatorname{FunctionLimit}(f,c,L)
  \;\Longleftrightarrow\;
  \operatorname{RightLimit}(f,c,L)\wedge\operatorname{LeftLimit}(f,c,L).
\]
\end{remark*}

\begin{remark*}[Negated quantified statement]
\[
\begin{aligned}
&\exists\varepsilon_0>0\;\forall\delta>0\;\exists x\in A\;
  0<|x-c|<\delta\;\wedge\; |f(x)-L|\geq\varepsilon_0 \\
&\quad\Longleftrightarrow
\Bigl[\exists\varepsilon_0>0\;\forall\delta>0\;\exists x\in A\;
  0<x-c<\delta\;\wedge\; |f(x)-L|\geq\varepsilon_0\Bigr] \\
&\qquad\vee
\Bigl[\exists\varepsilon_0>0\;\forall\delta>0\;\exists x\in A\;
  0<c-x<\delta\;\wedge\; |f(x)-L|\geq\varepsilon_0\Bigr].
\end{aligned}
\]
\end{remark*}

\begin{remark*}[Negation predicate reading]
\[
  \neg\,\operatorname{FunctionLimit}(f,c,L)
  \;\Longleftrightarrow\;
  \neg\,\operatorname{RightLimit}(f,c,L)\vee\neg\,\operatorname{LeftLimit}(f,c,L).
\]
\end{remark*}

\begin{remark*}[Failure modes]
A two-sided limit fails precisely when the right-hand limit fails, the
left-hand limit fails, or the two one-sided behaviours do not meet at the
same value.
\end{remark*}

\begin{remark*}[Failure modes]
\[
  \neg\,\operatorname{FunctionLimit}(f,c,L)
  \;\Longleftrightarrow\;
  \neg\,\operatorname{RightLimit}(f,c,L)\vee\neg\,\operatorname{LeftLimit}(f,c,L).
\]
\end{remark*}

\begin{remark*}[Interpretation]
A two-sided limit is exactly the agreement of the two one-sided tests.
Both sides must approach the same value.
\end{remark*}

\begin{dependencies}
\begin{itemize}
  \item \hyperref[def:limit-function]{Limit of a Function}
  \item \hyperref[def:right-hand-limit]{Right-Hand Limit}
  \item \hyperref[def:left-hand-limit]{Left-Hand Limit}
\end{itemize}
\end{dependencies}
